\subsection{\texorpdfstring{The homomorphism \(\bigoplus_{i} \mathrm{Hom}_{\mathrm{A}}(\mathrm{M}, \mathrm{N}_{i}) \longrightarrow \mathrm{Hom}_{\mathrm{A}}(\mathrm{M}, \bigoplus_{i} \mathrm{N}_{i})\)}{The homomorphism \textbackslash bigoplus\_\{i\} \textbackslash mathrm\{Hom\}\_\{\textbackslash mathrm\{A\}\}(\textbackslash mathrm\{M\}, \textbackslash mathrm\{N\}\_\{i\}) \textbackslash longrightarrow \textbackslash mathrm\{Hom\}\_\{\textbackslash mathrm\{A\}\}(\textbackslash mathrm\{M\}, \textbackslash bigoplus\_\{i\} \textbackslash mathrm\{N\}\_\{i\})}}

Let A be a ring, M an A-module, and \((\mathrm{N}_{i})_{i\in\mathrm{I}}\) a family of A-modules. With any element \((u_{i})\) of \(\bigoplus_{i}\operatorname{Hom}_{\mathrm{A}}(\mathrm{M},\mathrm{N}_{i})\) , we associate the element \(m\mapsto(u_{i}(m))\) of \(\operatorname{Hom}_{\mathrm{A}}(\mathrm{M},\bigoplus_{i}\mathrm{N}_{i})\) . We thus define a canonical homomorphism

\[
\varphi : \bigoplus_ {i} \operatorname{Hom} _ {\mathrm{A}} (\mathrm{M}, \mathrm{N} _ {i}) \longrightarrow \operatorname{Hom} _ {\mathrm{A}} \left(\mathrm{M}, \bigoplus_ {i} \mathrm{N} _ {i}\right).
\]

It is clear that \(\varphi\) is injective. Let u be an element of \(\operatorname{Hom}_{\mathrm{A}}(\mathrm{M},\bigoplus_{i}\mathrm{N}_{i})\) . Then u belongs to the image of \(\varphi\) if and only if the set of indices i such that \(pr_{i}\circ u\neq0\) is finite. This condition is automatically satisfied when the module M is finitely generated.

Consequently, if the module M is finitely generated, then the homomorphism \(\varphi\) is bijective.

